% TOP_PROBLEM: {"id": 2732, "title": "Kirby Problem 1.73", "category_id": 7, "category": {"id": 7, "name": "topology", "display_name": "Topology", "description": "Properties preserved under continuous deformations.", "slug": "topology", "order_index": 7, "created_at": "2026-07-31T15:26:25.671Z"}, "status": "open", "problem_number": "KP-1.73", "set_id": 11}
% FIRST_POSTED: 2026-10-01
% ENTRY_KIND: statement_only
% UnsolvedMath ID 2732; source catalogue code KP-1.73
% !TeX program = lualatex
% Definition and sources only; this file is not an LLM solution attempt.
\documentclass[11pt,a4paper]{article}
\usepackage[margin=25mm]{geometry}
\usepackage{fontspec}
\setmainfont{lmroman10-regular.otf}[BoldFont=lmroman10-bold.otf,
 ItalicFont=lmroman10-italic.otf,BoldItalicFont=lmroman10-bolditalic.otf]
\usepackage{amsmath,amssymb,mathtools}
\usepackage{unicode-math}
\setmathfont{latinmodern-math.otf}
\AtBeginDocument{\let\setminus\smallsetminus}
\usepackage{xurl,hyperref}
\hypersetup{colorlinks=true,urlcolor=blue}
\setcounter{secnumdepth}{0}
\setlength{\parindent}{0pt}
\setlength{\parskip}{5pt}
\setlength{\emergencystretch}{3em}
\begin{document}
\section*{Kirby Problem 1.73}\label{2732}
{\small UnsolvedMath ID 2732 · KP-1.73 · Topology · Kirby's Problems in Low-Dimensional Topology}
\subsection{Definitions and mathematical statement}
A pearl necklace is an embedded equilateral polygon in $\mathbb R^3$, together with equal closed balls centered at its vertices, whose diameter equals the polygon's edge length and whose interiors are pairwise disjoint. Adjacent balls are tangent because adjacent vertices are one edge length apart; additional tangencies are permitted. Knot type refers to the polygonal centerline, compared by ambient isotopy in $\mathbb R^3$.

By scaling, take the edge length to be one. Thus an $n$-pearl realization consists of cyclically ordered vertices $v_1,\ldots,v_n$ satisfying
\[
|v_{i+1}-v_i|=1,\qquad |v_i-v_j|\ge1\quad(i\ne j),
\]
with $v_{n+1}=v_1$, whose joining edges form an embedded polygon. Define $p(K)$ as the least possible $n$ for a necklace of knot type $K$.

Is $p(3_1)=15$, where $3_1$ is the trefoil knot? Either chirality gives the same value, since reflection preserves every metric constraint. The assertion combines existence of a fifteen-vertex realization with the impossibility of any realization using fourteen or fewer vertices.

The vertices can lie anywhere in Euclidean space; no lattice condition is imposed. Equal edge lengths alone are insufficient, because the balls at nonadjacent vertices must also avoid interior overlap. Likewise, a lower bound for polygons in one selected lattice does not bound all admissible pearl necklaces. The problem concerns this discrete geometric invariant, rather than the length of a flexible tube around a smooth knot.
\subsection{Short English statement}
Is fifteen the smallest number of equal nonoverlapping pearls whose polygonal centerline can form a trefoil?
\subsection{Sources}
\begin{itemize}
\item[S1] ULAM.AI and the UnsolvedMath Contributors, supplied problems.json, record 2732 (KP-1.73), Kirby's Problems in Low-Dimensional Topology. Catalogue identity and source transcription; CC BY 4.0.
\url{https://huggingface.co/datasets/ulamai/UnsolvedMath}.
\item[S2] K3: A New Problem List in Low-Dimensional Topology, Problem 1.73. Expanded formulation of the supplied catalogue statement.
\url{https://math.berkeley.edu/sites/default/files/surv-295-ruberman-watermarked-author-pdf.pdf}.
\end{itemize}
\end{document}
